% GENERATED FILE. Edit canonical lessons, then run npm run generate:books.
% canonical-source-hash: fnv1a64:4af2c072c981071d
% canonical-lessons: TA-C295-tuppu, TA-C295-tatayam, TA-C295-canteki, TA-C295-kurram, TA-C295-kurravali

\chapter{Clues and Suspects}
\label{ch:ta-clues-and-suspects}
\hlchaptermodality{295}

\begin{chapteropening}
\textbf{By the end of this chapter:} \emph{I can talk about clues, suspects and crimes.}

Complete the last lesson of chapter 295: I can talk about clues, suspects and crimes.
\end{chapteropening}

\section[\texorpdfstring{\ta{துப்பு} (tuppu)}{துப்பு (tuppu)}]{\ta{துப்பு} (tuppu) — a clue; a lead}

\label{lesson:TA-C295-tuppu}



\begin{quote}
Before the new one: say the Tamil for a robbery, then the Tamil for to happen.
\end{quote}

\subsection*{You'll want to know: \ta{துப்பு}}
\textbf{\ta{துப்பு}} — \emph{tuppu} — \textquotedblleft{}a clue; a lead\textquotedblright{}.

\begin{quote}
\ta{ஒரு} \ta{துப்பு} \ta{கிடைத்திருந்தால்}, \ta{காவலர்கள்} \ta{அவனைப்} \ta{பிடித்திருப்பார்கள்}.
\end{quote}

\begin{quote}
\emph{oru tuppu kiṭaittiruntāl, kāvalarkaḷ avaṉaip piṭittiruppārkaḷ.}
\end{quote}

\begin{quote}
If a clue had turned up, the police would have caught him.
\end{quote}

The finder takes -\ta{க்கு}: \ta{அவர்களுக்குத்} \ta{துப்பு} \ta{கிடைத்தது}, they found a clue.

\begin{cousinweb}
\textbf{\ta{துப்பு}} (\emph{tuppu}), \textquotedblleft{}a clue, a lead\textquotedblright{}, is native Dravidian: DEDR 3347 gives \ta{துப்பு}, \textquotedblleft{}spying, a trace, evidence of a crime\textquotedblright{}, under \ta{துருவு}, \textquotedblleft{}to search out, trace\textquotedblright{}.
\end{cousinweb}

\begin{grammarlens}[title={What you've built}]
One more word for telling what happened.
\end{grammarlens}

\subsection*{Your turn}
\begin{itemize}
\raggedright
  \item \emph{Say it:} \emph{tuppu}
  \item \emph{Say it:} \emph{tuppu}, once more
  \item \emph{Say it:} \emph{oru tuppu kiṭaittiruntāl, kāvalarkaḷ avaṉaip piṭittiruppārkaḷ.} — If a clue had turned up, the police would have caught him.
  \item \emph{Recall:} say the Tamil for a robbery, then the Tamil for to happen, then say \emph{tuppu} again
\end{itemize}

\begin{tcolorbox}[breakable,colback=teal!4,colframe=teal!35!black,title={Before you move on}]
What does \ta{துப்பு} mean? (\textquotedblleft{}A clue; a lead\textquotedblright{}.) Say it once more.
\end{tcolorbox}
\section[\texorpdfstring{\ta{தடயம்} (taṭayam)}{தடயம் (taṭayam)}]{\ta{தடயம்} (taṭayam) — a trace; a mark left behind}

\label{lesson:TA-C295-tatayam}



\begin{quote}
Before the new one: say the Tamil for to happen, then the Tamil for a clue.
\end{quote}

\subsection*{You'll want to know: \ta{தடயம்}}
\textbf{\ta{தடயம்}} — \emph{taṭayam} — \textquotedblleft{}a trace; a mark left behind\textquotedblright{}.

\begin{quote}
\ta{திருடன்} \ta{எந்தத்} \ta{தடயமும்} \ta{விட்டுச்} \ta{செல்லவில்லை}.
\end{quote}

\begin{quote}
\emph{tiruṭaṉ entat taṭayamum viṭṭuc cellavillai.}
\end{quote}

\begin{quote}
The thief left no trace behind.
\end{quote}

With -\ta{உம்} and a negative verb it means not a single one: \ta{எந்தத்} \ta{தடயமும்} \ta{இல்லை}.

\begin{cousinweb}
\textbf{\ta{தடயம்}} (\emph{taṭayam}), \textquotedblleft{}a trace, a mark left behind\textquotedblright{}, is a Tamil word whose origin is not settled in the sources at hand.
\end{cousinweb}

\begin{grammarlens}[title={What you've built}]
One more word for telling what happened.
\end{grammarlens}

\subsection*{Your turn}
\begin{itemize}
\raggedright
  \item \emph{Say it:} \emph{taṭayam}
  \item \emph{Say it:} \emph{taṭayam}, once more
  \item \emph{Say it:} \emph{tiruṭaṉ entat taṭayamum viṭṭuc cellavillai.} — The thief left no trace behind.
  \item \emph{Recall:} say the Tamil for to happen, then the Tamil for a clue, then say \emph{taṭayam} again
\end{itemize}

\begin{tcolorbox}[breakable,colback=teal!4,colframe=teal!35!black,title={Before you move on}]
What does \ta{தடயம்} mean? (\textquotedblleft{}A trace; a mark left behind\textquotedblright{}.) Say it once more.
\end{tcolorbox}
\section[\texorpdfstring{\ta{சந்தேகி} (cantēki)}{சந்தேகி (cantēki)}]{\ta{சந்தேகி} (cantēki) — to suspect}

\label{lesson:TA-C295-canteki}



\begin{quote}
Before the new one: say the Tamil for a clue, then the Tamil for a trace.
\end{quote}

\subsection*{You'll want to know: \ta{சந்தேகி}}
\textbf{\ta{சந்தேகி}} — \emph{cantēki} — \textquotedblleft{}to suspect\textquotedblright{}.

\begin{quote}
\ta{காவலர்கள்} \ta{அவனுடைய} \ta{நண்பனைச்} \ta{சந்தேகித்தார்கள்}.
\end{quote}

\begin{quote}
\emph{kāvalarkaḷ avaṉuṭaiya naṇpaṉaic cantēkittārkaḷ.}
\end{quote}

\begin{quote}
The police suspected his friend.
\end{quote}

Past (I …): \ta{சந்தேகித்தேன்}, \emph{cantēkittēṉ}.

The one suspected takes -\ta{ஐ}; the noun is \ta{சந்தேகம்}, a doubt.

\begin{cousinweb}
\textbf{\ta{சந்தேகி}} (\emph{cantēki}) is a Tamil verb built on \ta{சந்தேகம்}, \textquotedblleft{}doubt\textquotedblright{}, which is borrowed from Sanskrit \emph{saṃdeha}, \textquotedblleft{}doubt, uncertainty\textquotedblright{} (Monier-Williams).
\end{cousinweb}

\begin{grammarlens}[title={What you've built}]
One more word for telling what happened.
\end{grammarlens}

\subsection*{Your turn}
\begin{itemize}
\raggedright
  \item \emph{Say it:} \emph{cantēki}
  \item \emph{Say it:} \emph{cantēki}, once more
  \item \emph{Say it:} \emph{kāvalarkaḷ avaṉuṭaiya naṇpaṉaic cantēkittārkaḷ.} — The police suspected his friend.
  \item \emph{Recall:} say the Tamil for a clue, then the Tamil for a trace, then say \emph{cantēki} again
\end{itemize}

\begin{tcolorbox}[breakable,colback=teal!4,colframe=teal!35!black,title={Before you move on}]
What does \ta{சந்தேகி} mean? (\textquotedblleft{}To suspect\textquotedblright{}.) Say it once more.
\end{tcolorbox}
\section[\texorpdfstring{\ta{குற்றம்} (kuṟṟam)}{குற்றம் (kuṟṟam)}]{\ta{குற்றம்} (kuṟṟam) — a crime; a fault}

\label{lesson:TA-C295-kurram}



\begin{quote}
Before the new one: say the Tamil for a trace, then the Tamil for to suspect.
\end{quote}

\subsection*{You'll want to know: \ta{குற்றம்}}
\textbf{\ta{குற்றம்}} — \emph{kuṟṟam} — \textquotedblleft{}a crime; a fault\textquotedblright{}.

\begin{quote}
\ta{திருடுவது} \ta{ஒரு} \ta{குற்றம்}.
\end{quote}

\begin{quote}
\emph{tiruṭuvatu oru kuṟṟam.}
\end{quote}

\begin{quote}
Stealing is a crime.
\end{quote}

A -\ta{அம்} noun: the plural is \ta{குற்றங்கள்}, crimes.

\begin{cousinweb}
\textbf{\ta{குற்றம்}} (\emph{kuṟṟam}) is native Dravidian: DEDR 1851 lists it with \ta{குறு}, \textquotedblleft{}short, defective\textquotedblright{}, so a crime is at root a fault, a falling short.
\end{cousinweb}

\begin{grammarlens}[title={What you've built}]
One more word for telling what happened.
\end{grammarlens}

\subsection*{Your turn}
\begin{itemize}
\raggedright
  \item \emph{Say it:} \emph{kuṟṟam}
  \item \emph{Say it:} \emph{kuṟṟam}, once more
  \item \emph{Say it:} \emph{tiruṭuvatu oru kuṟṟam.} — Stealing is a crime.
  \item \emph{Recall:} say the Tamil for a trace, then the Tamil for to suspect, then say \emph{kuṟṟam} again
\end{itemize}

\begin{tcolorbox}[breakable,colback=teal!4,colframe=teal!35!black,title={Before you move on}]
What does \ta{குற்றம்} mean? (\textquotedblleft{}A crime; a fault\textquotedblright{}.) Say it once more.
\end{tcolorbox}
\section[\texorpdfstring{\ta{குற்றவாளி} (kuṟṟavāḷi)}{குற்றவாளி (kuṟṟavāḷi)}]{\ta{குற்றவாளி} (kuṟṟavāḷi) — a culprit; a guilty person}

\label{lesson:TA-C295-kurravali}



\begin{quote}
Before the new one: say the Tamil for to suspect, then the Tamil for a crime.
\end{quote}

\subsection*{You'll want to know: \ta{குற்றவாளி}}
\textbf{\ta{குற்றவாளி}} — \emph{kuṟṟavāḷi} — \textquotedblleft{}a culprit; a guilty person\textquotedblright{}.

\begin{quote}
\ta{நீதிபதி} \ta{அவனைக்} \ta{குற்றவாளி} \ta{என்று} \ta{சொன்னார்}.
\end{quote}

\begin{quote}
\emph{nītipati avaṉaik kuṟṟavāḷi eṉṟu coṉṉār.}
\end{quote}

\begin{quote}
The judge said he was the culprit.
\end{quote}

-\ta{ஆளி} marks a person: \ta{குற்றம்} + \ta{ஆளி}, one who has the fault.

\begin{cousinweb}
\textbf{\ta{குற்றவாளி}} (\emph{kuṟṟavāḷi}) is a Tamil compound built from \ta{குற்றம்}, \textquotedblleft{}fault, crime\textquotedblright{} (DEDR 1851), and -\ta{ஆளி}, \textquotedblleft{}a person who has or does\textquotedblright{}.
\end{cousinweb}

\begin{grammarlens}[title={What you've built}]
One more word for telling what happened.
\end{grammarlens}

\subsection*{Your turn}
\begin{itemize}
\raggedright
  \item \emph{Say it:} \emph{kuṟṟavāḷi}
  \item \emph{Say it:} \emph{kuṟṟavāḷi}, once more
  \item \emph{Say it:} \emph{nītipati avaṉaik kuṟṟavāḷi eṉṟu coṉṉār.} — The judge said he was the culprit.
  \item \emph{Recall:} say the Tamil for to suspect, then the Tamil for a crime, then say \emph{kuṟṟavāḷi} again
\end{itemize}

\begin{tcolorbox}[breakable,colback=teal!4,colframe=teal!35!black,title={Before you move on}]
What does \ta{குற்றவாளி} mean? (\textquotedblleft{}A culprit; a guilty person\textquotedblright{}.) Say it once more.
\end{tcolorbox}
